\section{Introduction}
\label{sec:introduction}

Language-model agents use tools to inspect schemas, execute Python or SQL, retrieve documentation, and produce analytical reports \citep{schick2023toolformer,yao2023react,qin2023toolllm}. Their answers depend on both the computation and how they interpret its output. Successful execution alone does not establish that the returned answer is correct.

Most tool-use evaluations emphasize tool selection, valid arguments, execution errors, or end-task success. Data workflows also encounter failures that preserve a normal-looking interface: an aggregate uses the wrong denominator, a statistic is attached to the wrong entity, or a retrieved data dictionary is stale. These errors can pass format checks and still change the answer. We call this setting \emph{silent tool poisoning}: the tool executes successfully, but the observation supplied to the agent contains controlled factual or semantic corruption.

ToxicBench asks whether agents detect, verify, or adopt this corrupted evidence. Each task pairs the same query, data, model, and adapter in clean and poisoned environments. Numerical operators alter magnitudes, signs, ratios, or rankings; semantic operators alter entity bindings, column meanings, metadata, or retrieved evidence. Alongside task success, trajectory metrics separate adoption without checking from adoption after verification and successful recovery. This distinction matters because an agent may request another computation yet still select the wrong conclusion.

We evaluate framework adapters for LangGraph, smolagents, DA-Agent, and AutoGen where feasible. The cross-model suite contains 34 numerical and 24 semantic/schema instances; the revised expanded GPT-only analysis retains 118 of 120 released instances after reference auditing. PandasAI's historical final-output corruption is reported separately because that agent cannot respond after the injected change. Across the observation-level settings, high clean-task success coexists with substantial poisoned-task degradation.

We compare Double-pass, Verification-only, and Generic Guard with equal route and token limits, followed by repeated poisoning over numerical, semantic/schema, and multi-table workflows. These protocols distinguish another attempt from an explicit request to verify the first answer. A 240-trajectory human audit with third-rater adjudication exposed answer-selection errors and informed the revised scorer. Its subsequent recheck is a development evaluation, not independent validation. Trace analysis separates corrupted follow-up checks from unmodified observations whose relevance still needs assessment.

Our contributions are:
\begin{itemize}
    \item A data-analysis benchmark with paired clean and poisoned runs, seven numerical operators, five semantic/schema operators, and a multi-table extension.
    \item A trajectory protocol separating adoption, checking, and recovery, with exposure counts, human-audit evidence, and versioned parser and reference-answer sensitivity analyses.
    \item Cross-framework and cross-model experiments, equal-budget verification controls, clustered sensitivity checks, and 1,596 recorded repeated-poison trajectories (1,572 retained after reference exclusions).
\end{itemize}
